% Section II (Amauri Martins): to be written. The mathematical framework refers to this
% section for the parameter box P of the cable and overhead-line inputs and for the
% per-unit-length series impedance z'_b(f; p) and shunt admittance y'_b(f; p).
% Findings of the IEEE 30-bus study (HarmonicPowerModels.jl pull request 35) to report here:
% - Inputs: register of 173 rows with a source and class for every nominal value and interval;
%   39 candidate inputs, one coordinate per physical family (e.g., one conductor temperature
%   for all 22 cables); a candidate is admitted when moving it alone to either end of its
%   interval changes some sequence z or y by at least 0.1 % at 50 or 650 Hz; 23 admitted
%   (cables 10, ACSR phase conductors 8, earth wire 4, earth resistivity 1, log scale,
%   10--1000 Ohm m). Largest single effects: cable conductor temperature 26 % on z0, earth
%   resistivity 14.6 % on z0, ACSR temperature 6.9 % on z1. Of the 16 dropped inputs, 9 have
%   no effect (e.g., insulation resistivity, XLPE loss tangent) and 7 change z or y by at most
%   0.036 %. The endpoint screening does not test interactions between coordinates.
% - Coupled inputs: each box point is mapped onto the physical set (phase temperature at least
%   ambient; wire-diameter deviations scaled back to the EN 50182 diameter (+-1 %) and mass
%   (+-2 %) limits); the maps are the identity on feasible points and continuous.
% - Uncertainty belongs on the physical inputs: boxing the flattened LCM equivalents (each
%   conductor as one tube, each insulation as layers) needs 35 coordinates for guaranteed
%   containment, against 23 physical inputs, as one physical input moves every size of its
%   family by a slightly different law (study in Amauri's fork).
% - Bonding and transposition: screens solidly bonded at both ends, earth wire grounded at each
%   tower, ideal transposition, no cross-bonding.
% - Installed length: one coordinate per link through the length of the exact pi-model
%   (supported by HarmonicPowerModels.jl; to be offered as an option in the study).
\section{Cable and Overhead-Line Parameter Uncertainty}
\label{sec:parameter_uncertainty}
% $\mathcal{P}$ and $\bar{p}$ are \ParamSet and \ParamVec of the mathematical framework; use the macros once they move to nomenclature.tex.

The per-unit-length series impedance~$\bar{z}'_{\Line}(\Frequency;\bar{p})$ and shunt admittance~$\bar{y}'_{\Line}(\Frequency;\bar{p})$ follow from physical inputs of the cables and overhead lines, accounting for the corresponding material and dimensional tolerances, and computed with the Julia package \textsc{LineCableModels.jl}~\cite{lcm_2026}. Each cable or line is a system of tubular conductors separated by insulating and semiconducting layers, in the air or in the earth~\cite{ametani_cable_2015}. For~$n$ conductors, the $n \times n$ matrices of series impedance and shunt admittance per unit length are
\begin{IEEEeqnarray}{ r C l }
    Z'(\Frequency;\bar{p}) &=& Z'_{\text{c}} + Z'_{\text{i}} + Z'_{\text{e}}, \label{eq:pul_impedance} \\
    Y'(\Frequency;\bar{p}) &=& \mathrm{j} 2\pi\Frequency \left( P'_{\text{i}} + P'_{\text{e}} \right)^{-1}, \label{eq:pul_admittance}
\end{IEEEeqnarray}
in which $Z'_{\text{c}}$ collects the surface impedances of the conductors, $Z'_{\text{i}}$ the impedances of the insulating layers and $Z'_{\text{e}}$ the earth-return impedances; and $P'_{\text{i}}$ and $P'_{\text{e}}$ are the potential coefficients of the insulating layers and of the surrounding air and earth.

% The surface impedances are ratios of modified Bessel functions of~$r\sqrt{\mathrm{j}2\pi\Frequency\mu/\rho}$ for a conductor of radius~$r$, permeability~$\mu$ and resistivity~$\rho$. The factor~$\sqrt{\pi\Frequency\mu/\rho}$ is the reciprocal of the skin depth, which in aluminum falls from 12.0\,mm at 50\,Hz to 3.3\,mm at 650\,Hz, below the 13.25\,mm radius of a 500\,mm$^{2}$ core. The earth return penetrates 712\,m at 50\,Hz and 197\,m at 650\,Hz into a 100\,$\Omega\cdot$m earth, so the zero-sequence parameters, whose currents return through screens, earth wires and earth, vary most. 
The effect of temperature~$\vartheta$, in $^{\circ}$C, on metal resistivity is described by 
\begin{equation}
    \rho(\vartheta) = \rho_{20} \left[ 1 + \kappa_{20} \left( \vartheta - 20 \right) \right],
    \label{eq:resistivity_temperature}
\end{equation}
where $\rho_{20}$ is the resistivity at 20\,$^{\circ}$C and $\kappa_{20}$ its temperature coefficient~\cite{iec62641_2022,iec60287_1_1_2006}. In overhead lines, temperature also sets the sag through the elastic thermal change-of-state equation, and each conductor lies at its support height less two thirds of the sag.

Cable screens are assumed bonded and solidly earthed at both cable terminals and the earth wires at every tower, with circuits ideally transposed. This yields circulant-reduced $3 \times 3$ matrices, expressed as zero-, positive- and negative-sequence parameters ~$\bar{z}'_{\Line}$ and~$\bar{y}'_{\Line}$ of Section~\ref{sec:parametric_network}.

Table~\ref{tab:parameter_box} summarizes the cables, lines and earth of the IEEE 30-bus case, whose cable sizes and ACSR codes follow from the branch currents, rated under IEEE~738~\cite{ieee738_2023}. The intervals are the acceptance limits of the standards cited in Table~\ref{tab:parameter_box}, operating and environmental ranges for the temperatures and the earth, and study intervals otherwise. Two continuous maps, equal to the identity at feasible points, keep every point of the box physical: the phase temperature is at least the ambient one, and the wire-diameter deviations of each ACSR code are scaled back onto the diameter and mass limits of EN~50182~\cite{en50182_2001}.

Physical inputs are represented as parameter space coordinates per input of a family, e.g., one conductor temperature for all 22 cables or, for a per-size quantity, it moves every size within its own interval. Of 39 candidate inputs, the 23 that change some sequence impedance or admittance by at least 0.1\,\% at 50 or 650\,Hz, when moved alone to an end of their interval, form the coordinate. The complete dataset, with the source and class of every value, and the specifications and case files are available in the HarmonicPowerModels.jl repository~\cite{hpm_2026}.

\begin{table*}[!tb]
    \centering
    \caption{Cables, overhead lines and earth of the IEEE 30-bus case: layout, and nominal value and interval of the 11 most influential of the 23 uncertain inputs, with the largest relative change of a sequence parameter at 50 or 650\,Hz when the input alone moves to an end of its interval; $z_{0}$, $z_{1}$ and $y_{0}$, $y_{1}$ are the zero- and positive-sequence series impedance and shunt admittance per unit length. }
    \label{tab:parameter_box}
    \footnotesize\setlength{\tabcolsep}{4pt}
    \begin{tabular}{ >{\raggedright\arraybackslash}p{0.42\textwidth} >{\raggedright\arraybackslash}p{0.20\textwidth} l r }
        quantity & nominal & interval & largest change \\ \hline
        \multicolumn{4}{l}{\textit{cables, 33\,kV, 22 links in touching trefoil, center 0.8\,m deep, screens bonded at both ends}} \\
        aluminum conductor cross section~\cite{iec60228_2023} & 400--800\,mm$^{2}$ & by link & -- \\
        conductor DC resistance at 20\,$^{\circ}$C~\cite{iec60228_2023} & 0.0367--0.0778\,$\Omega$/km & $-8$\,\% to nominal & 4.4\,\% $z_{1}$ \\
        conductor diameter~\cite{iec60228_2023} & 23.0--34.2\,mm & nominal to $+0.4$\,mm & 1.3\,\% $y_{0}$ \\
        XLPE insulation thickness~\cite{iec60502_2_2014} & 8.0\,mm & 7.6 to 8.4\,mm & 4.4\,\% $y_{0}$ \\
        XLPE relative permittivity~\cite{iec60287_1_1_2006} & 2.4 & 2.1 to 2.5 & 12.5\,\% $y_{1}$ \\
        screen DC resistance at 20\,$^{\circ}$C, 48 copper wires of 1.35\,mm & 0.263\,$\Omega$/km & 0.251 to 0.263\,$\Omega$/km & 3.8\,\% $z_{0}$ \\
        jacket thickness~\cite{iec60502_2_2014} & 2.85--3.30\,mm & nominal to $+35$\,\% & 2.0\,\% $z_{1}$ \\
        conductor temperature~\cite{iec60502_2_2014} & 20\,$^{\circ}$C & 10 to 90\,$^{\circ}$C & 26\,\% $z_{0}$ \\
        \multicolumn{4}{l}{\textit{overhead lines, 132\,kV, 12 single-circuit lines in delta, span 300\,m}} \\
        ACSR phase conductors, 30 aluminum and 7 steel wires~\cite{en50182_2001} & Lion, Goat, Bear, Panther & by line & -- \\
        phase positions $(x, y)$ at the tower & $(\pm 2.5, 18.0)$, $(0, 22.33)$\,m & -- & -- \\
        earth wire, 7 steel wires of 3.15\,mm~\cite{iec63248_2022}, position & $(0, 29.0)$\,m & -- & -- \\
        reference sag at 20\,$^{\circ}$C, phases and earth wire & 6.0 and 5.4\,m & -- & -- \\
        phase temperature & 20\,$^{\circ}$C & $-5$ to 75\,$^{\circ}$C & 6.9\,\% $z_{1}$ \\
        ambient and earth-wire temperature & 20\,$^{\circ}$C & $-5$ to 45\,$^{\circ}$C & 3.2\,\% $z_{1}$ \\
        earth-wire steel relative permeability & 100 & 50 to 300 & 3.1\,\% $z_{0}$ \\
        \multicolumn{4}{l}{\textit{earth, homogeneous}} \\
        resistivity, logarithmic scale & 100\,$\Omega\cdot$m & 10 to 1000\,$\Omega\cdot$m & 14.6\,\% $z_{0}$ \\
        relative permittivity & 10 & -- & -- \\
        \multicolumn{4}{l}{\textit{12 further inputs}} \\
        semiconducting screens; wire diameters and lay ratios~\cite{en50182_2001,iec62641_2022,iec63248_2022}; metal resistivities; phase steel permeability & & & $\leq 0.62$\,\%
    \end{tabular}
\end{table*}
